% GENERATED FILE. Edit canonical lessons, then run npm run generate:books.
% canonical-source-hash: fnv1a64:7c2b2e11c315e512
% canonical-lessons: JA-C186-amaru, JA-C186-chokin, JA-C186-tameru, JA-C186-ryoukin, JA-C186-zeikin

\chapter{Savings, Fees and Tax}
\label{ch:ja-savings-fees-and-tax}
\hlchaptermodality{186}

\begin{chapteropening}
\textbf{By the end of this chapter:} \emph{I can talk about savings, fees and tax.}

Complete the last lesson of chapter 186: I can talk about savings, fees and tax.
\end{chapteropening}

\section[\texorpdfstring{\ja{あまる} (amaru)}{あまる (amaru)}]{\ja{あまる} (amaru) — to be in surplus; to be left over}

\label{lesson:JA-C186-amaru}



\begin{quote}
Before the new one: say the Japanese for to cost, then the Japanese for to be enough.
\end{quote}

\subsection*{You'll want to know: \ja{あまる}}
\textbf{\ja{あまる}} — \emph{amaru} — \textquotedblleft{}to be in surplus; to be left over\textquotedblright{}.

\begin{quote}
\textbf{\ja{おかねが} \ja{すこし} \ja{あまります。}} — \emph{okane ga sukoshi amarimasu} — \textquotedblleft{}A little money is left over.\textquotedblright{}
\end{quote}

Polite \textbf{\ja{あまります}}, \emph{amarimasu}; a godan verb.

\begin{cousinweb}
\textbf{\ja{あまる}} (\emph{amaru}) is a native Japanese word (wago), inherited rather than borrowed from Chinese; its kanji was attached later as a kun'yomi (JMdict, KANJIDIC2).
\end{cousinweb}

\begin{grammarlens}[title={What you've built}]
One more word for everyday life at A2.
\end{grammarlens}

\subsection*{Your turn}
\begin{itemize}
\raggedright
  \item \emph{Say it:} \emph{amaru}
  \item \emph{Say it:} \emph{amaru}, once more
  \item \emph{Say it:} \emph{okane ga sukoshi amarimasu} — A little money is left over.
  \item \emph{Recall:} say the Japanese for to cost, then the Japanese for to be enough, then say \emph{amaru} again
\end{itemize}

\begin{tcolorbox}[breakable,colback=teal!4,colframe=teal!35!black,title={Before you move on}]
What does \ja{あまる} mean? (\textquotedblleft{}To be in surplus; to be left over\textquotedblright{}.) Say it once more.
\end{tcolorbox}
\section[\texorpdfstring{\ja{ちょきん} (chokin)}{ちょきん (chokin)}]{\ja{ちょきん} (chokin) — savings; money put aside}

\label{lesson:JA-C186-chokin}



\begin{quote}
Before the new one: say the Japanese for to be enough, then the Japanese for to be in surplus.
\end{quote}

\subsection*{You'll want to know: \ja{ちょきん}}
\textbf{\ja{ちょきん}} — \emph{chokin} — \textquotedblleft{}savings; money put aside\textquotedblright{}.

\begin{quote}
\textbf{\ja{ぎんこうに} \ja{ちょきんが} \ja{あります。}} — \emph{ginkō ni chokin ga arimasu} — \textquotedblleft{}I have savings in the bank.\textquotedblright{}
\end{quote}

\begin{cousinweb}
\textbf{\ja{ちょきん}} (\emph{chokin}) is Sino-Japanese (kango): two Chinese characters read with their on'yomi, readings borrowed from Chinese. They are \emph{cho}, \textquotedblleft{}store up\textquotedblright{} (Middle Chinese \emph{trjoX}) and \emph{kin}, \textquotedblleft{}gold, money\textquotedblright{} (\emph{kim}) (KANJIDIC2; Middle Chinese in Baxter's transcription).
\end{cousinweb}

\begin{grammarlens}[title={What you've built}]
One more word for everyday life at A2.
\end{grammarlens}

\subsection*{Your turn}
\begin{itemize}
\raggedright
  \item \emph{Say it:} \emph{chokin}
  \item \emph{Say it:} \emph{chokin}, once more
  \item \emph{Say it:} \emph{ginkō ni chokin ga arimasu} — I have savings in the bank.
  \item \emph{Recall:} say the Japanese for to be enough, then the Japanese for to be in surplus, then say \emph{chokin} again
\end{itemize}

\begin{tcolorbox}[breakable,colback=teal!4,colframe=teal!35!black,title={Before you move on}]
What does \ja{ちょきん} mean? (\textquotedblleft{}Savings; money put aside\textquotedblright{}.) Say it once more.
\end{tcolorbox}
\section[\texorpdfstring{\ja{ためる} (tameru)}{ためる (tameru)}]{\ja{ためる} (tameru) — to save up; to store, to pile up}

\label{lesson:JA-C186-tameru}



\begin{quote}
Before the new one: say the Japanese for to be in surplus, then the Japanese for savings.
\end{quote}

\subsection*{You'll want to know: \ja{ためる}}
\textbf{\ja{ためる}} — \emph{tameru} — \textquotedblleft{}to save up; to store, to pile up\textquotedblright{}.

\begin{quote}
\textbf{\ja{いもうとは} \ja{おかねを} \ja{ためます。}} — \emph{imōto wa okane o tamemasu} — \textquotedblleft{}My younger sister saves up money.\textquotedblright{}
\end{quote}

Polite \textbf{\ja{ためます}}, \emph{tamemasu}; an ichidan verb.

\begin{cousinweb}
\textbf{\ja{ためる}} (\emph{tameru}) is a native Japanese word (wago), inherited rather than borrowed from Chinese; its kanji was attached later as a kun'yomi (JMdict, KANJIDIC2).
\end{cousinweb}

\begin{grammarlens}[title={What you've built}]
One more word for everyday life at A2.
\end{grammarlens}

\subsection*{Your turn}
\begin{itemize}
\raggedright
  \item \emph{Say it:} \emph{tameru}
  \item \emph{Say it:} \emph{tameru}, once more
  \item \emph{Say it:} \emph{imōto wa okane o tamemasu} — My younger sister saves up money.
  \item \emph{Recall:} say the Japanese for to be in surplus, then the Japanese for savings, then say \emph{tameru} again
\end{itemize}

\begin{tcolorbox}[breakable,colback=teal!4,colframe=teal!35!black,title={Before you move on}]
What does \ja{ためる} mean? (\textquotedblleft{}To save up; to store, to pile up\textquotedblright{}.) Say it once more.
\end{tcolorbox}
\section[\texorpdfstring{\ja{りょうきん} (ryōkin)}{りょうきん (ryōkin)}]{\ja{りょうきん} (ryōkin) — fee; charge, fare}

\label{lesson:JA-C186-ryoukin}



\begin{quote}
Before the new one: say the Japanese for savings, then the Japanese for to save up.
\end{quote}

\subsection*{You'll want to know: \ja{りょうきん}}
\textbf{\ja{りょうきん}} — \emph{ryōkin} — \textquotedblleft{}fee; charge, fare\textquotedblright{}.

\begin{quote}
\textbf{\ja{りょうきんは} \ja{いくらですか。}} — \emph{ryōkin wa ikura desu ka} — \textquotedblleft{}How much is the fee?\textquotedblright{}
\end{quote}

\begin{cousinweb}
\textbf{\ja{りょうきん}} (\emph{ryōkin}) is Sino-Japanese (kango): two Chinese characters read with their on'yomi, readings borrowed from Chinese. They are \emph{ryō}, \textquotedblleft{}fee, materials\textquotedblright{} (Middle Chinese \emph{lew}) and \emph{kin}, \textquotedblleft{}gold, money\textquotedblright{} (\emph{kim}) (KANJIDIC2; Middle Chinese in Baxter's transcription).
\end{cousinweb}

\begin{grammarlens}[title={What you've built}]
One more word for everyday life at A2.
\end{grammarlens}

\subsection*{Your turn}
\begin{itemize}
\raggedright
  \item \emph{Say it:} \emph{ryōkin}
  \item \emph{Say it:} \emph{ryōkin}, once more
  \item \emph{Say it:} \emph{ryōkin wa ikura desu ka} — How much is the fee?
  \item \emph{Recall:} say the Japanese for savings, then the Japanese for to save up, then say \emph{ryōkin} again
\end{itemize}

\begin{tcolorbox}[breakable,colback=teal!4,colframe=teal!35!black,title={Before you move on}]
What does \ja{りょうきん} mean? (\textquotedblleft{}Fee; charge, fare\textquotedblright{}.) Say it once more.
\end{tcolorbox}
\section[\texorpdfstring{\ja{ぜいきん} (zeikin)}{ぜいきん (zeikin)}]{\ja{ぜいきん} (zeikin) — tax}

\label{lesson:JA-C186-zeikin}



\begin{quote}
Before the new one: say the Japanese for to save up, then the Japanese for fee.
\end{quote}

\subsection*{You'll want to know: \ja{ぜいきん}}
\textbf{\ja{ぜいきん}} — \emph{zeikin} — \textquotedblleft{}tax\textquotedblright{}.

\begin{quote}
\textbf{\ja{ぜいきんを} \ja{はらいます。}} — \emph{zeikin o haraimasu} — \textquotedblleft{}I pay tax.\textquotedblright{}
\end{quote}

\begin{cousinweb}
\textbf{\ja{ぜいきん}} (\emph{zeikin}) is Sino-Japanese (kango): two Chinese characters read with their on'yomi, readings borrowed from Chinese. They are \emph{zei}, \textquotedblleft{}tax, duty\textquotedblright{} (Middle Chinese \emph{sywejH}) and \emph{kin}, \textquotedblleft{}gold, money\textquotedblright{} (\emph{kim}) (KANJIDIC2; Middle Chinese in Baxter's transcription).
\end{cousinweb}

\begin{grammarlens}[title={What you've built}]
One more word for everyday life at A2.
\end{grammarlens}

\subsection*{Your turn}
\begin{itemize}
\raggedright
  \item \emph{Say it:} \emph{zeikin}
  \item \emph{Say it:} \emph{zeikin}, once more
  \item \emph{Say it:} \emph{zeikin o haraimasu} — I pay tax.
  \item \emph{Recall:} say the Japanese for to save up, then the Japanese for fee, then say \emph{zeikin} again
\end{itemize}

\begin{tcolorbox}[breakable,colback=teal!4,colframe=teal!35!black,title={Before you move on}]
What does \ja{ぜいきん} mean? (\textquotedblleft{}Tax\textquotedblright{}.) Say it once more.
\end{tcolorbox}
